% Въ лѣто Господне ҂вк҃ѕ.
%
% Господи Іисусе Христе, Сыне Божій, помилуй меня грѣшнаго.

\documentclass[20pt]{extarticle}

\usepackage[a5paper,margin=.6cm,top=0pt,bottom=0pt]{geometry}

%\usepackage{pgfpages}
%\pgfpagesuselayout{2 on 1}[a4paper, landscape]

\usepackage{fontspec}
\usepackage{xkeyval}
\usepackage{microtype}
\usepackage{parskip}

\usepackage{polyglossia}
\setmainlanguage{churchslavonic}

\usepackage{churchslavonic}
\usepackage{lettrine}
\usepackage{graphicx}

\usepackage{hyperref}
\hypersetup{
	pdftitle={На лiтiи стiхиры},
	pdfauthor={Православная Церковь}
}


\makeatletter

\def\@texttop{\vfill}
\def\@textbottom{\vfill}

\makeatother


\newcommand{\cuL}[1]{\scalebox{1.5}{\cuKinovarColor#1}}
\newcommand{\cul}[1]{\mbox{\cuKinovarColor#1}}

\renewcommand\star[1][\relax]{\smash{\rlap{\scalebox{1.5}{%
	\hspace{-1pt}\raisebox{-1.5pt}{%
	\if\relax\detokenize{#1}\relax%
		\cuKinovarColor%
	\fi%
	*%
}}}}}

\newfontfamily\churchslavonicfont[Script=Cyrillic,Ligatures=TeX,HyphenChar=_]{Acathist}
%\newfontfamily\cyrillicfontsf{DejaVu Sans}
%\newfontfamily\noto{Noto Sans Symbols 2}

\hyphenpenalty=10000
\setlength\emergencystretch{4em}
\def\hyph{{\_}{}{}\penalty-10000}

%\linespread{1.2}
%\setlength{\parindent}{1cm}
\setlength{\parskip}{0pt}

\renewcommand{\title}[1]{{\centering\Large\cuKinovar{#1}\par}}
\newcommand\sub[1]{{\centering\normalsize\cuKinovar{#1}\par}}
\renewcommand{\emph}[1]{\mbox{\textls[200]{#1}}}


\begin{document}

\pagestyle{empty}
\large


\sub{Тропа́рь, гла́съ д҃.}
\vspace{4pt}

\begin{center}

\cuL{Р}жⷭ҇тво̀ твоѐ, бцⷣе дв҃о\star, ра́дость возвѣстѝ все́й вселе́ннѣй\star:
и҆зъ тебѐ бо возсїѧ̀ со́лнце пра́вды, хрⷭ҇то́съ бг҃ъ на́шъ\star, \\ и҆ разрꙋши́въ клѧ́твꙋ,
дадѐ благослове́нїе\star[], \\ и҆, ѹ҆праздни́въ сме́рть, дарова̀ на́мъ живо́тъ вѣ́чный.

\end{center}


\vspace{12pt}
\sub{Конда́къ, гла́съ д҃.}

\begin{center}

\cuL{І҆}ѡакі́мъ и҆ а҆́нна поноше́нїѧ безча́дства\star, и҆ а҆да́мъ и҆ є҆́ва ѿ тлѝ сме́ртныѧ
свободи́стасѧ\star, пречⷭ҇таѧ\star, во ст҃ѣ́мъ ржⷭ҇твѣ̀ твое́мъ\star. То пра́зднꙋютъ
и҆ лю́дїе твоѝ, вины̀ прегрѣше́нїй изба́вльшесѧ\star, внегда̀ зва́ти тѝ\star[]:
непло́ды ражда́етъ бцⷣꙋ и҆ пита́тельницꙋ жи́зни на́шеѧ.

\end{center}


\vspace{12pt}
\sub{Велича́нїе:}

\begin{center}

\cuL{В}елича́емъ тѧ̀\star, прест҃а́ѧ дв҃о\star, и҆ чти́мъ ст҃ы́хъ твои́хъ
роди́телей\star, и҆ всесла́вное сла́вимъ\star[] ржⷭ҇тво̀ твоѐ.

\end{center}
\vspace{12pt}

\end{document}
